\chapter{Introduction}

TODO: Chapter on why NIMEs for real-time musical performance are useful and exciting?\\


Recent developments in new interfaces for musical expression (NIMEs) which can be applied in real-time musical performance have broadly fallen into two categories. Higher level symbolic approaches break down musical harmony and structure symbolically, breaking down the complexities of musical expression into discrete components which can be logically reasoned with. By converting the expression of a human performer of the instrument into this symbolic language, a NIME can then dynamically interact with the performer by performing calculations in this logical space (include citation for a paper like this?). These symbolic methods are excellent for generating structured, coherent, and theoretically sensical musical outputs. However, this approach suffers from the severe restrictions it places on the domain of musical expression, which often leads to robotic and less creative sounding results (need citation). In contrast, other methods have used recurrent neural networks (RNNs) trained on human performance on an instrument to repeatedly predict the most likely next step in a performance. When used for NIMEs, these techniques solve many of the issues of symbolic approaches, as they are formulated to continuously predict the most human sounding continuation of a musical performance. However, they suffer in longer term structure, as RNNs are focused on repeated short term thinking, and hence suffer the same pitfall as most greedy algorithms. These approaches will repeatedly generate great human sounding creative musical phrases, but these musical phrases usually have no coherent long term structure, and it is clear that the AI is trying to continuously make the most impressive five second segment of a performance over and over again, instead of making a coherent and impressive longer performance (cite An Interactive Musical Prediction System with Mixture Density Recurrent Neural Networks at some point).\\


This current state of the art in intelligent musical instruments presents multiple important research questions which this paper aims to answer. Firstly, how can the strengths of symbolic approaches and RNN-based methods in NIMEs be combined into a single system which to achieve both short-term creativity and long-term structural coherence in real-time musical performances? Secondly, in terms of implementation, what are the most effective structural parameters for creating long-term coherence, and how can these be quantified and combined into a heuristic? Thirdly, how would a combined prediction system compare to existing approaches in terms of musical quality, structural coherence, and real-time performance capabilities, as assessed by human performers? By addressing these questions, we aim to develop a new framework for interactive musical prediction systems (IMPS) which combine both symbolic and neural network based approaches, to create better NIMEs as assessed by human performers.\\


Our combined approach for developing real-time interactive NIMEs uses an RNN based IMPS as a baseline algorithm, and extends it with a symbolic structural heuristic. As a baseline algorithm we use the mixture density recurrent neural network (MDRNN) IMPS developed by (cite An Interactive Musical Prediction System with Mixture Density Recurrent Neural Networks), which uses an RNN based approach. We then modify the prediction system. At each point where a prediction must be made, instead of outputting the first prediction made by the RNN, the RNN is run multiple times recursively, generating an expanding tree of possible branches the performance could continue along while sounding human like. On each of these branches, we then measure a set of structural parameters. Examples of these parameters are rhythmic uniformity, harmonic consonance, and pitch variation. Branches which are more consistent with the previous state of these structural parameters are given higher heuristic values, leading the choice of next prediction to favour musical directions which don't break from the long term structural norm in an incoherent way. When interacting with a human performer, these structural parameters, and the rate of change of each, will be tracked and used to inform the AI when it takes over from the human in the performance, ensuring that it continues the performance in the same structural direction that the human performer was moving in. This also allows for sudden breaks from structural parameters in places where it is deemed to be a reasonable direction for the piece, for example the beginning of a bridge section. We also present an interface for the human performer to control the structural parameters of the AI while it plays in a live setting, and for new and unique structural parameters to be trained from human performance, allowing for longer term parameters which are difficult to symbolically code, such as the mood of a piece, to be encoded using machine learning instead.\\


Using RNN predictions to determine on a short time scale which musical directions are appropriate ensures that the performance will sound more human and creative than symbolic approaches. The addition of symbolically encoded structural parameters and a tree search based system of choosing the musical direction of a piece negates the disadvantages of RNN based approaches, as it provides a framework for coherent long term structure, and balances out the greedy next step prediction process of a purely RNN based approach. (briefly explain the structure of the paper and what each section will cover here).